\section{Holography and the Lower Faces}

The committed substrate is four-dimensional and cannot be pictured whole, but the same family has lower-dimensional members that can, and they are useful for more than intuition. Some properties are easiest to confirm cleanly on a simpler face, and holography is the clearest example. The two-dimensional $\{7,3\}$ heptagrid and the three-dimensional $\{5,3,4\}$ honeycomb are not the thesis, they are the faces on which a hard measurement becomes tractable.

The holographic principle says that the information in a region of space is captured by its boundary, and that the entanglement of a boundary region is fixed by the geometry of a surface reaching into the bulk. This is the Ryu-Takayanagi law, and on the $\{7,3\}$ face it is not assumed, it is measured. The entanglement of a boundary interval grows as the logarithm of the interval's length, the signature of a one-dimensional boundary anchoring a geodesic in a negatively curved bulk, against a flat control where the same quantity grows linearly. The discriminating control could have failed and did not, so the holographic law is a genuine emergent property of the dynamics on the crystal (depth L3).

\begin{result}[Area, not volume]
The ground state of the emergent field obeys an area law, its entanglement scaling with the boundary of a region rather than its volume, while a thermal state crosses over to a volume law. Read as a black hole, a region's entropy scales with the area of its horizon and not the volume it encloses, the Bekenstein-Hawking relation, recovered on the substrate (depth L3).
\end{result}

The negative curvature that makes holography work also makes the bulk a natural error-correcting code, and on the $\{5,3,4\}$ face this can be built explicitly. Tiling the bulk with the perfect tensors of a holographic code gives a structure in which a bit written deep in the interior is protected against erasure at the boundary, with the protection growing as the code reaches deeper, the distance rising as a power of the depth. A bit spread across the bulk survives a bounded erasure that destroys a bit stored in one place. This is the same robustness that the later chapters lean on when they read the depth of the bulk as a memory, and it is built and verified here on the curved face before it is put to that further use (depth L2 to L3). The lesson to carry forward is that the substrate is holographic by its geometry, so that depth into the bulk is at once a scale, a protection, and, eventually, an inwardness.
